% r03_fill_K6: proof of (K_0), part IX: the existence of the axis-lattice limits, the forward bound, and the assembly.

\subsubsection*{The existence of the walled axis-lattice limits}

\begin{theorem}[the walled torus limit on the axis lattice]\label{rot:thm:H6W}
For every walled configuration $\fc=(J,\psi,n)$ of the class of \Cref{rot:thm:CT} and every real smooth $P_J$-periodic $2$-form $G$ the limit
\[
\Wv_\fc(tG)\defeq\lim_{k\to\infty}\frac{\Zun_{\Xk,T_1}(tG,h_0)}{\Zun_{\Xk,T_1}(0,h_0)}=\lim_{k\to\infty}\langle e^{tF(G)}\rangle_{\fc,\Xk}
\]
exists for every real $t$; it is the same number in the free-boundary convention of \Cref{rot:def:tori} and in King's torus convention at the shifted constant source $h_0-\Delta_T(\eps_k)/2$ (\Cref{rot:prop:CT}(c)), and it equals the limit of the torus-convention ratios $\Zun^{\rm tor}_{\Xk,T_1}(tG,h_\infty)/\Zun^{\rm tor}_{\Xk,T_1}(0,h_\infty)$ at the frozen wall $h_\infty\defeq h_0-\Delta_T(0)/2$. The convergence is locally uniform in $t\in\C$, the limit is entire, and all moments of $F(G)$ converge.
\end{theorem}

This is King's torus theorem \cite[Thm.~2.1]{K1} with the wall as his scalar source, read on the normalised ratio, in the only form in which it is true at a wall: the convergence of King's effective action \cite[Thm.~3.4]{K1} modulo a field-independent constant of its scalar-source sector, which cancels in the ratio, with the drift-free stability bound \cite[(2.23)]{K1} replaced by a polynomial envelope. The proof is given in the following three groups of statements; it compares the walled ratio, which is King's torus-convention ratio at a wall shifted by a bounded constant that moves with $k$ by $O(\eps_k)$, with the torus-convention ratio at a frozen wall on the same lattice, proves the frozen family Cauchy by telescoping one step at a time, and bounds the effect of the constant shift of the source by a per-site bound with its volume factor against the rate of \Cref{rot:prop:CT}(b).

% r03_h6w_a: proof of Theorem rot:thm:H6W, part 1: conventions and the bridge, the a priori bounds, King's flow on the axis lattice, the two one-lattice stability inputs, and the reduction to the small-field ratios.
\providecommand{\Ztor}{\ensuremath{\Zun^{\rm tor}}}
\providecommand{\hinf}{\ensuremath{h_\infty}}

The proof of \Cref{rot:thm:H6W} occupies the rest of this paragraph, up to the subsubsection on the forward bound. It is a one-lattice statement: only the axis lattice $\Xk=\Lk$ on the torus $T\defeq T_1=\DJ$ occurs, with King's standard block-spin flow at block factor $5$ and King's own corner-rooted blocks (the centred blocks of \Cref{rot:conv:C} play no role here), and neither the rotated lattice, nor the hybrid flow, nor \Cref{rot:thm:sandwich} is used. Throughout, $\fc=(J,\psi,n)$ is in the class of \Cref{rot:thm:CT}, $h_0\defeq-n\psi$, $G$ is a real smooth $P_J$-periodic $2$-form sampled in the frame of $\Xk$, $t\in\R$, and $r=r(k)$, $\delta_r=5^r\eps_k=5^{-\lfloor\vr k\rfloor}$ are the depth and the coarse spacing of \Cref{rot:lem:r}. The argument compares two ratios on the same lattice: the walled curvature ratio, which by the bridge of \Cref{rot:prop:CT} is King's torus-convention ratio at a wall that moves with $k$ by a constant of order $\eps_k$, and the torus-convention ratio at a frozen wall; the frozen family is Cauchy by King's renormalisation theorem read modulo the field-independent constant of its scalar-source sector, and the moving family stays close to the frozen one because a constant shift of the scalar source changes the effective action by a bounded amount per site, which the rate $|u_k|\le C\eps_k$ of \Cref{rot:prop:CT}(b) beats.

\subsubsection*{Conventions, the bridge and the a priori bounds}

\begin{lemma}[the moving and the frozen wall]\label{rot:lem:H6conv}
Write $\Ztor_k(t;h)\defeq\Ztor_{\Xk,T}(tG,h)$ for the torus-convention functional of \Cref{rot:prop:CT}(c) and
\[
h_k\defeq h_0-\tfrac12\Delta_T(\eps_k)\mathbf 1,\qquad\hinf\defeq h_0-\tfrac12\Delta_T(0)\mathbf 1,\qquad u_k\defeq h_k-\hinf=\tfrac12\big(\Delta_T(0)-\Delta_T(\eps_k)\big),
\]
a constant with $|u_k|\le C\eps_k$. Then for every $k$ and every real $t$
\[
r_k(t)\defeq\langle e^{tF(G)}\rangle_{\fc,\Xk}=\frac{\Zun_{\Xk,T}(tG,h_0)}{\Zun_{\Xk,T}(0,h_0)}=\frac{\Ztor_k(t;h_k)}{\Ztor_k(0;h_k)},\qquad s_k(t)\defeq\frac{\Ztor_k(t;\hinf)}{\Ztor_k(0;\hinf)} .
\]
The family $\{h_k\}_{k}\cup\{\hinf\}$ is bounded in every $C^N(T)$ norm, $\|h\|_{C^N}\le\|h_0\|_{C^N}+\Delta_*/2$; $\Xk/P_J$ is King's lattice on his torus \cite[(2.21)]{K1}, so the right side of the identity for $r_k$ is literally the object of \cite[Thms.~3.1, 3.4]{K1} with the scalar source $h_k$, and King's constants, which depend on the sources only through bounds on finitely many derivatives (\cite[pp.~664--665, 677]{K1}: ``the sources $g,h$ are bounded by $C$'', ``their derivatives are uniformly bounded''), are uniform over the family. The physical mass $m^2$ is the same in both conventions and is never changed.
\end{lemma}
\begin{proof}
The first identity is the walled-ratio display of \Cref{rot:def:Zun} at $g=0$, with $-w=h_0$ (see also \Cref{rot:lem:walledalg}(c)): the wall is the scalar source $h_0$ up to a constant that cancels in the ratio. The second is \Cref{rot:prop:CT}(c) at $\Lambda=\Xk$, whose prefactor does not depend on $G$ and cancels in the ratio. The existence of $\Delta_T(0)=\lim_{s\to0}\Delta_T(s)$ and $|u_k|\le C\eps_k$ are the rate clause of \Cref{rot:prop:CT}(b), and $|\Delta_T(\eps_k)|,|\Delta_T(0)|\le\Delta_*$ give the $C^N$ bound, the shifts being constants. No sign condition on the scalar source enters King's bounds: $\lambda>0$ makes every single-site measure finite for every real source.
\end{proof}

\begin{lemma}[a priori bounds on the ratios]\label{rot:lem:H6A}
Let $C_G\defeq\sup_kq_{\Xk}(G)/2<\infty$ (\Cref{rot:lem:q}). For every $k$, every real smooth $P_J$-periodic $h$ and every real $t$,
\[
1\le\frac{\Ztor_k(t;h)}{\Ztor_k(0;h)}\le e^{t^2q_{\Xk}(G)/2}\le e^{C_Gt^2},\qquad\Big|\frac{\Ztor_k(z;h)}{\Ztor_k(0;h)}\Big|\le\frac{\Ztor_k(\Re z;h)}{\Ztor_k(0;h)}\le e^{C_G(\Re z)^2}\quad(z\in\C),
\]
and the same in the free-boundary convention. Consequently, if a sequence of such ratios converges for every real $t$, it converges locally uniformly on $\C$ to an entire function, and all moments of $F(G)$ converge to the Taylor coefficients of the limit.
\end{lemma}
\begin{proof}
The torus-convention walled state is a BFS measure of King type in the sense of the paragraph after \Cref{rot:def:Zun}: the torus counterterm and Wick constant only change the Gaussian factor of the single-site measures and a constant. By \Cref{lem:s02:GD} with $F_{\Xk}(G)=\ip A{\ell_{\Xk}}$ (\Cref{rot:lem:CM}), $\langle e^{tF(G)}\rangle\le e^{t^2q_{\Xk}(G)/2}$; by Jensen's inequality and $\langle F(G)\rangle=0$ under the charge conjugation of \Cref{rot:def:tori}, $\langle e^{tF(G)}\rangle\ge1$. For complex $z$, $|e^{zF}|=e^{(\Re z)F}$. The functions $z\mapsto\Ztor_k(z;h)/\Ztor_k(0;h)$ are entire (the integral converges absolutely for every complex $z$ by the same Gaussian bound) and are bounded on every disc $|z|\le\mathsf R$ by $e^{C_G\mathsf R^2}$, uniformly in $k$; convergence on $\R$ then gives locally uniform convergence on $\C$ by Vitali's theorem \cite{Hor90}, and Cauchy's estimates give the moments.
\end{proof}

\subsubsection*{King's flow on the axis lattice}

For a lattice $X_{k'}$ and a depth $r'$ with $5^{r'}\eps_{k'}<1$, King's standard flow \cite[\S3]{K1} on $X_{k'}/P_J$ after $r'$ steps produces the unit lattice $T^{(r')}$ of physical spacing $\delta\defeq5^{r'}\eps_{k'}$, the small-field indicator $\chi_{r'}$ \cite[(3.2)]{K1}, on whose support $|\phi_{r'}|\le p(\delta)(\lambda\delta)^{-1/4}$, and the effective action $S^{(r')}_{k'}(tG,h)$ of \Cref{rot:def:eff}, here in the torus-convention reading for which \cite{K1} is written: its Gaussian parts (quadratic forms, $\ln N_{r'}$, $\ln Z_{r'}(A^{(r')})$) are defined at the mass $m^2$ without sources \cite[(2.13)--(2.16)]{K1}, and the sources enter only through the curvature vertex \cite[(3.20)]{K1} and the $h$-vertex \cite[(3.21)]{K1} of the perturbative part $E'_{r'}$, expanded to King's order $\bar n=13$ with a remainder. Put
\[
I^{(r')}_{k'}(t;h)\defeq\int\chi_{r'}\,e^{-S^{(r')}_{k'}(tG,h)}\,dA_{r'}\,d\phi_{r'},\qquad
\rho^{(r')}_{k'}(t;h)\defeq\frac{I^{(r')}_{k'}(t;h)}{I^{(r')}_{k'}(0;h)},
\]
and $I_k\defeq I^{(r(k))}_k$, $\rho_k\defeq\rho^{(r(k))}_k$ at the lattice's own depth. Two facts are used throughout: the unit lattice has $|T^{(r')}|=|T|\delta^{-3}$ sites (\cite[(2.15)]{K1} in unit-lattice counting); and $5^r\eps_k=5^{r+1}\eps_{k+1}=\delta_r$, so that $\Xk$ at depth $r$ and $X_{k+1}$ at depth $r+1$ live on the same unit lattice $T^{(r)}$ with the same $\chi_r$ and the same field bounds (\cite[p.~656]{K1}: ``since $L^k\eps=L^{k+n}\eps'$, the fields $A_k,\phi_k$ have the same bounds in each integral'').

\subsubsection*{The two stability inputs at the walls}

\begin{proposition}[King's upper bound on the axis lattice, at a wall]\label{rot:prop:H6S1}
For $h$ in the family of \Cref{rot:lem:H6conv}, every $k$ and every admissible depth $r'$, with $\delta=5^{r'}\eps_k$,
\[
e^{-C\delta^\sigma|T|}\,I^{(r')}_k(t;h)\le\Ztor_k(t;h)\le e^{C\delta^\sigma|T|}\,I^{(r')}_k(t;h)+B_k,\qquad B_k\defeq e^{-p(\delta)^2+C|T|}\eps_k^{-q},
\]
with $\sigma>0$, $C=C(t,G,h_0,|T|)$ locally uniform in $t$, and $q=q(t,G,h_0)<\infty$, all independent of $k$, of $r'$ and of $h$ in the family.
\end{proposition}
\begin{proof}
This is \cite[Thm.~3.1, (3.3)--(3.4)]{K1} for King's torus with the sources $(tG,h)$, with the additive large-field term in the form in which King uses it (\cite[(3.11)]{K1}, \cite[(4.7)]{K2}), at block factor $5$ on the torus $\DJ$ aligned with King's blocks: \Cref{thm:a06:BALA}(iii) with \Cref{hyp:k01:BALint}, in the torus convention for which it is printed, so that no bridge is crossed. Two readings are made, the same two as in \Cref{rot:prop:H3}: the uniformity over the bounded source family (\Cref{rot:lem:H6conv}); and the envelope. The factor $e^{C|T|}$ of the large-field term is the ultraviolet stability bound \cite[(2.23)]{K1}, which fails as printed when $\int h\ne0$ (\Cref{rot:rem:log}); what the inductive proof supplies at a wall is a bounded $h$-linear $c$-number per step, $e^{C|T|+C'(1+k)|\int h|}=e^{C|T|}\eps_k^{-q}$. The loss is harmless against $e^{-p(\delta_r)^2}$, since $p(\delta_r)^2$ grows at least quadratically in $k$ (\Cref{rot:lem:r}) while $\ln\eps_k^{-q}=qk\ln5$ is linear. The sources enter the proof through the vertices (3.20)--(3.21) only, with constants depending on $\sup|G|$, $\sup|\nabla G|$, $\sup|h|$, $\sup|\nabla h|$ (\cite[pp.~664--665, 677]{K1}).
\end{proof}

\begin{proposition}[the polynomial lower bound at the wall, torus convention]\label{rot:prop:H6S2}
For $h$ in the family of \Cref{rot:lem:H6conv} and every $k$: $\Ztor_k(0;h)\ge e^{-C_1|T|}\eps_k^{q_1}$, with $C_1$, $q_1$ independent of $k$ and of $h$ in the family.
\end{proposition}
\begin{proof}
This is \Cref{rot:thm:FS} on $\Lambda=\Xk$, $T=T_1$, for the bounded family ($\|h\|_\infty\le H\defeq n\|\psi\|_\infty+\Delta_*/2$; the sign of $h$ is not used there), read in the torus convention: the proof of \Cref{rot:thm:FS} moves the wall into the bare scalar coefficient, compares with a higher King point by Jensen's inequality and the monotonicity of $\langle\sum_x\eps^3|\phi|^2\rangle$ in that coefficient, and uses the source-free stability bounds \cite[(1.14)]{Ba1}, \cite[(1.3)]{Ba2} at the higher points together with the integrated counterterm identity and the perturbative bounds on the one-point Wick function. Each of these steps is available for the torus-convention model with the same constants: \cite[(1.14)]{Ba1} and \cite[(1.3)]{Ba2} are printed for it (\Cref{thm:a06:BALA}(iii)); its bare coefficient, Wick constant and vacuum constant differ from the free-boundary ones by $\Delta_T(s)$, $\Delta c_T(s)$, $\Delta E_T(s)$, which are bounded uniformly in $s$ (\Cref{rot:prop:CT}(b)) and shift the King points by bounded amounts only; and the perturbative coefficients of the torus one-point function differ from the infinite-lattice ones by image sums, exponentially small in $\ell(T)$. The bound reads $\ln\Ztor_k(0;h)\ge-C|T|-C(1+k)|T|-qk\ln5$, which is the claim with $q_1=q+C|T|/\ln5$.

At the mass-increasing wall the same conclusion also follows from Jensen's inequality as in \Cref{rot:prop:J}, in the torus convention: $\Ztor_k(0;h)\ge\Ztor_k(0;0)\exp[\sS_{\Xk}(h)\,w^{\rm tor}_k]$ with $w^{\rm tor}_k\defeq\langle|\phi(x)|^2\rangle^{\rm tor}_{0,0}-2c_T(\eps_k)$, independent of $x$ by translation invariance ($P_J\subset\Xk$), where $\Ztor_k(0;0)\ge e^{-C|T|}$ is \cite[(1.14)]{Ba1} and $\sS_{\Xk}(h)=-n\sS_{\Xk}(\psi)-\frac12\Delta_T|T|<0$ by the wall condition of \Cref{rot:thm:CT}; together with an upper bound $w^{\rm tor}_k\le C(1+k)$ on the source-free one-point Wick function, which is not established here and not needed, this gives the claim with $q_1=C|\sS_{\Xk}(h)|/\ln5$. This is the only place in the proof where the sign of the wall would enter; the route through \Cref{rot:thm:FS} does not use it.
\end{proof}

\subsubsection*{The ratios are small-field ratios up to $o(1)$}

\begin{lemma}[the large-field term is negligible]\label{rot:lem:H6stepA}
Let $h$ be in the family of \Cref{rot:lem:H6conv}, $r'$ an admissible depth for $\Xk$, $\delta=5^{r'}\eps_k$, $x\defeq C\delta^\sigma|T|$ with the constant of \Cref{rot:prop:H6S1}, and $b_k(h)\defeq2B_k/\Ztor_k(0;h)$. Then
\[
b_k(h)\le2e^{-p(\delta)^2+(C+C_1)|T|}\eps_k^{-q-q_1},
\]
and whenever $b_k\le\frac12$ and $2x\le\ln2$:
\begin{equation}\label{rot:eq:H6stepA}
e^{-2x}(1-b_k)\,\rho^{(r')}_k(t;h)\le\frac{\Ztor_k(t;h)}{\Ztor_k(0;h)}\le e^{2x}\rho^{(r')}_k(t;h)+b_k,\qquad
\rho^{(r')}_k(t;h)\le4e^{C_Gt^2},\qquad
\Big|\frac{\Ztor_k(t;h)}{\Ztor_k(0;h)}-\rho^{(r')}_k(t;h)\Big|\le4e^{C_Gt^2}(e^{2x}-1)+(1+4e^{C_Gt^2})\,b_k .
\end{equation}
At the lattice's own depth $r=r(k)$, $x=C\delta_r^\sigma|T|\to0$ and $b_k(h)\to0$ as $k\to\infty$ (\Cref{rot:lem:r}: $p(\delta_r)^2$ grows at least quadratically in $k$, the exponent of $\eps_k^{-q-q_1}$ linearly), so $|r_k(t)-\rho_k(t;h_k)|\to0$ and $|s_k(t)-\rho_k(t;\hinf)|\to0$ for every real $t$, locally uniformly in $t$.
\end{lemma}
\begin{proof}
Write $Z(t)\defeq\Ztor_k(t;h)$, $I(t)\defeq I^{(r')}_k(t;h)$, $\rho\defeq\rho^{(r')}_k(t;h)$, $B\defeq B_k$, $b\defeq b_k(h)$. The bound on $b$ is \Cref{rot:prop:H6S1,rot:prop:H6S2}. By \Cref{rot:prop:H6S1} at $t$ and at $0$: $Z(t)\le e^xI(t)+B$ and $Z(0)\ge e^{-x}I(0)$ give $Z(t)/Z(0)\le e^{2x}\rho+B/Z(0)\le e^{2x}\rho+b$; and $Z(t)\ge e^{-x}I(t)=e^{-x}\rho I(0)$ with $I(0)\ge e^{-x}(Z(0)-B)$ give $Z(t)/Z(0)\ge e^{-2x}\rho(1-B/Z(0))\ge e^{-2x}(1-b)\rho$. Hence, by \Cref{rot:lem:H6A}, $\rho\le e^{2x}(1-b)^{-1}Z(t)/Z(0)\le2e^{2x}e^{C_Gt^2}\le4e^{C_Gt^2}$. For the difference: $Z(t)/Z(0)-\rho\le(e^{2x}-1)\rho+b$ and $\rho-Z(t)/Z(0)\le\rho[(1-e^{-2x})+b]\le\rho[(e^{2x}-1)+b]$; the factor $(1+4e^{C_Gt^2})$ on $b$ covers both sides. The last statement is the case $r'=r(k)$, $h\in\{h_k,\hinf\}$, with the identities of \Cref{rot:lem:H6conv}.
\end{proof}

% r03_h6w_b: proof of Theorem rot:thm:H6W, part 2: graphs and the census of the source subgraphs, the restricted form of King's renormalisation theorem, and the Cauchy estimate at the frozen wall modulo the field-free constant.

\subsubsection*{Graphs, degrees and the census of the source subgraphs}

King's perturbative part $P^{(r)}_k(tG,h)$ is a finite sum of connected renormalised graphs $H$ \cite[(3.29), (3.37)]{K1} with the vertices \cite[(3.18)--(3.26)]{K1}; the number of graphs depends on $\bar n$ only, each graph is multilinear in its source insertions, and its value on $\supp\chi_r$ is $\sum_{\{y\},\{z\}}E^{(r)}(H;\{y\},\{z\})\times(\text{external fields at }\{y\})$, with a kernel that depends on the coarse background $A^{(r)}$ through its lines $G_r(A^{(r)})$ \cite[(3.17)]{K1}, so that a graph without external legs still depends on the background. Degrees are King's \cite[(3.30)--(3.31)]{K1}, as in \Cref{rot:lem:S1}, computed on the point-split graph (the two summation points of a block product as vertices, as in \Cref{rot:cor:OF}), with two conventions. The curvature vertex \cite[(3.20)]{K1} is counted with $D_v=\frac52$, the lattice curl of its photon leg having been summed by parts onto the smooth $G$ (exact on the torus, $|d^*_\eta G|\le C\|\nabla G\|_\infty$); with the curl charged to the vertex, $D_v=\frac32$ and the tree with two curvature vertices would have degree $0$, so the summation by parts is necessary. The $h$-vertex \cite[(3.21)]{K1} has $D_v=2$ (\Cref{rot:lem:S1}). Every connected subgraph containing an $h$-vertex and a curvature vertex therefore has $D\ge2+\frac52-3=\frac32>0$, and the smallest vacuum graph with two curvature vertices and two current vertices has $D=5$.

Each graph value is split at zero coarse fields, $E(H)=E(H)|_0+[E(H)-E(H)|_0]$, where $|_0$ means $(A_r,\phi_r)=(0,0)$, and we put $E_N(H)\defeq E(H)$ if $H$ has an external scalar leg or a curvature vertex and $E_N(H)\defeq E(H)-E(H)|_0$ otherwise: the field- or source-dependent part. The field-free, $G$-free values $E(H)|_0$ of the remaining graphs ($h$-linear or $h$-nonlinear vacuum graphs without curvature vertex, and the $h$-free vacuum energy) are numbers depending on $k$, $r$, $h$ only; nothing is claimed about their convergence, and for the order-$e^2$ $h$-linear family it fails (\Cref{rot:rem:log}: the Wick constant drifts by $c_*\ln5$ per step).

\begin{lemma}[the census of the source subgraphs]\label{rot:lem:H6census}
Expand every graph in the coarse background as in \cite[(3.43)--(3.55)]{K1} (\Cref{rot:lem:CB}(1) for one flow): $E(H)$ becomes a finite sum of graphs $\tilde H$ with lines at zero background, explicit background legs $\tilde A$ at block points, and remainder factors, and $E(H)|_0$ is the core $\tilde H_0$ without background legs. Then:
\begin{enumerate}[label=\textup{(\alph*)},leftmargin=*]
\item Every $\tilde H$ entering $E_N(H)$ is connected and has an external leg (scalar, background or $dG$) or a curvature vertex.
\item For such $\tilde H$, every connected ordered subgraph $\gamma$ (lines with scale indices below $r$) has strictly positive King degree or vanishes identically. In detail: if $\gamma$ contains a block point, $D(\gamma)\ge1$ ($\ge2$ with a source vertex); if $\gamma$ contains an $h$-vertex, $D(\gamma)\ge\frac12$ unless $\gamma$ is a vacuum graph with exactly one $h$-vertex and no curvature vertex (then $\gamma=\tilde H$, excluded by (a)) or $\gamma$ is the $h$-loop through a single current vertex \cite[(3.22)]{K1} with its photon leg external ($D=0$: the current tadpole of \Cref{rot:lem:S1} with the $h$-insertion), which vanishes identically at zero background ($\Tr q=0$); if $\gamma$ contains a curvature vertex and no $h$-vertex, $D(\gamma)\ge\frac12$ unless $\gamma$ is the curvature vertex joined by a photon line to a current tadpole ($D=0$), which vanishes identically at zero background for the same reason; if $\gamma$ is source-free, $D(\gamma)>0$ after Ba{\l}aban's resummed rewriting \cite[(3.5)]{Ba3}, the input of \Cref{rot:cor:OF}.
\item The rewriting of the source-free subdivergences inside a graph that carries source vertices (Ba{\l}aban's Taylor subtraction at the subdivergence with its counterterm partner, \cite[(3.5)]{Ba3}, \cite[(3.83)--(3.85)]{K1}) does not lower the degree of a prefix component that contains a source vertex below the value counted in (b).
\item The only connected subgraphs containing a source vertex with King degree $\le0$ in the raw count are the field-free $h$-linear vacuum graphs of order $\le2$ in the couplings, the family $\cT_h$ of \Cref{rot:lem:S1} (which have no external leg, hence occur only as whole graphs), and the two vanishing exceptions of (b).
\end{enumerate}
Consequently the minimum $\alpha_0$ of the degrees of the non-vanishing ordered subgraphs over the finitely many graphs and orderings is strictly positive.
\end{lemma}
\begin{proof}
(a) The core $\tilde H_0$ is the same in $E(H)$ and $E(H)|_0$ and cancels in $E_N(H)$ when $H$ has no external scalar leg and no curvature vertex; every other $\tilde H$ carries a background leg, an external $dG$ line or a remainder factor (the remainders are bounded separately, \Cref{rot:lem:CB}(2)). (b) The degree arithmetic is that of \Cref{rot:lem:S1}: inserting an $h$-vertex into a scalar line adds $2$; source-free connected subgraphs with an external leg have $D\ge-1$ except the current tadpole ($D=-2$), so a subgraph with an $h$-vertex and an external leg has $D\ge1$ except the $h$-loop through one current vertex ($D=0$); a vacuum subgraph with exactly one $h$-vertex has $D\le0$ only at order $\le2$ in the couplings, the exhaustive enumeration of \Cref{rot:lem:S1}, and it is then the whole graph; with two or more $h$-vertices $D\ge1$. For a curvature vertex, which attaches by a photon leg and contributes $\frac52$, the same enumeration leaves the one exception named. Both exceptions vanish at zero background because the current one-point function of the loop vanishes, $\Tr q=0$, exactly as for the current tadpole in \Cref{rot:lem:S1}. A component containing a block point has degree $\ge1$ by the point-split count of \Cref{rot:cor:OF} (every component containing a block point has Ba{\l}aban's degree $\ge1$, \cite[(2.1)--(2.2)]{Ba3}), and a source vertex adds at least $1$ more. (c) Ba{\l}aban's Taylor subtraction at a source-free subdivergence $S$ raises the degree of $S$ by $\theta_S$ and lowers by $\theta_S$ the degree of the external line of $S$ that receives the derivative; the construction, run on the whole graph with the source vertices counted as vertices, rewrites $S$ only in orderings in which every line of $S$ is finer than every external line of $S$. Let $\gamma$ be a component of an ordered subgraph. If the line that receives the derivative lies in $\gamma$, then every line of $S$, being finer, lies in the same ordered subgraph, and $S$, connected to that line, lies in $\gamma$: the two changes cancel. Otherwise the degree of $\gamma$ can only increase. By induction over the forest of subdivergences, finest first, no component falls below its value in (b); the source vertices, which are never inside a rewritten $S$, keep their contributions $2$ and $\frac52$. (d) is read off the same arithmetic.
\end{proof}

\subsubsection*{The restricted renormalisation theorem}

King's renormalisation theorem \cite[Thm.~3.5]{K1} asserts the per-site bound (3.38) and the convergence (3.39) of the graph kernels for all graphs. As printed it is false on the order-$e^2$ $h$-linear field-free family (\Cref{rot:rem:log}), and the convergence of the field-free, $G$-free constants is not needed here. What is used is the following restricted form.

\begin{proposition}[the restricted renormalisation theorem]\label{rot:prop:H6R}
Let $h$ be in the family of \Cref{rot:lem:H6conv}, $t$ real, and let $\gamma_*,\beta_*,\sigma$ be the exponents of \cite[Thm.~3.4]{K1}.
\begin{enumerate}[label=\textup{(\roman*)},leftmargin=*]
\item \textup{(convergence)} For every graph $H$ and every $k$, with $r=r(k)$, the field- or source-dependent parts obey, after the sum over the positions, for any $\gamma'<\gamma_*$, $\beta'>\frac72$, $\sigma'<\sigma$,
\[
\sum_{\text{positions}}\big|E^{(r+1)}_{N,k+1}(H)-E^{(r)}_{N,k}(H)\big|\times|\text{external fields}|\le C\big[5^{-\gamma'r}\delta_r^{-\beta'}+\delta_r^{\sigma'}\big]|T|\quad\text{on }\supp\chi_r,
\]
the two kernels living on the common unit lattice $T^{(r)}$ of $\Xk$ at depth $r$ and $X_{k+1}$ at depth $r+1$; $C=C(t,G,h_0,|T|)$ is independent of $k$ and locally uniform in $t$.
\item \textup{(per-site bound)} For every graph $H$ with at least one $h$-vertex, other than the degree-$0$ members of the order-$e^2$ $h$-linear field-free family,
\[
|E^{(r)}_k(H;\{y\},\{z\})|\le C(1+k-r)^{p_0}\,\delta_r^{\varrho_H}\,p(\delta_r)^{n_1}\,e^{-\delta\dist(\{y\},\{z\})},\qquad\varrho_H\ge0,
\]
and for the degree-$0$ members of that family the same bound with the extra factor $(1+r)$, which is necessary there. For a graph at non-zero background the bound is read through the expansion of \Cref{rot:lem:H6census}: the core at zero background falls under the second clause if $H$ is in that family and under the first otherwise, and every other term has an external leg and falls under the first. $C$ is linear in $\sup|h|$ for a distinguished $h$-vertex and depends on the other sources through sup norms; $p_0$, $n_1$ depend on $\bar n$ only.
\end{enumerate}
\end{proposition}
\begin{proof}
King's proof of (3.39) \cite[pp.~660--662]{K1} reduces it to the convergence of a single graph with one line replaced by the difference of the two flows' lines (3.56), and that to \cite[Prop.~3.6]{K1}, the one-flow estimate for a renormalised family with one element weakened by a small degree, which King asserts on p.~665 ``by redoing the analysis of Sect.~3.4'' and which is \Cref{rot:cor:OF}. \Cref{rot:cor:OF} applies to a family exactly when every connected ordered subgraph of its non-vanishing graphs has strictly positive degree after Ba{\l}aban's rewriting: a zero-degree ultraviolet subgraph containing the replaced line would give $\sum_j(5^j\eta)^{-\gamma}\sim5^{\gamma r}$ and cancel the gain $5^{-\gamma r}$. By \Cref{rot:lem:H6census} this holds, with the minimum $\alpha_0>0$, for every $\tilde H$ entering $E_N(H)$. The replaced line is a fine line of \cite[(3.17)]{K1} with scale index below $r$; the lines with indices between $r$ and $k$ occur only on the counterterm lines of the $\dm$ graphs, never at a source vertex, so the counterterm-attachment clause \cite[(3.78)]{K1} and the one-vertex attachment are the source-free statements and are unchanged by the presence of the sources. The remaining steps are King's: the kernel bounds \cite[Thm.~3.3, Props.~3.7--3.9]{K1} for the lines and the fine-field reconstructions, printed for the torus convention at block factor $5$ (\Cref{thm:a06:BALA}); the sampling of the sources $x'\to x$ (p.~665), which costs $C\eps_k\|\nabla h\|_\infty+C\eps_k\|\nabla G\|_\infty\le C5^{-\gamma r}$; the tree-level $h$-term $\delta_r^2\sum\eta^3h|\phi^{(r)}|^2$, which has no internal line and is compared directly through the convergence of the fine-field reconstructions $a_rG_rQ_r^*$ \cite[Thm.~3.3, \S3.6]{K1}; the non-polynomial terms of (3.14) by \cite[\S3.6, (3.89)--(3.98), Prop.~3.10]{K1}; the Taylor remainders beyond order $\bar n$, $C\delta_r^\sigma|T|$ per flow (p.~665); and the position sum (3.40)--(3.42), which gives $|T|$ with one free position, the others tied by the exponential tree decay. The factor $(1+k-r)^{p_0}$ of \Cref{rot:cor:OF} is a power of $\log\delta_r^{-1}$ and is absorbed by $\gamma'<\gamma_*$, $\sigma'<\sigma$, and the field norms $(\lambda\delta_r)^{-s/4}p(\delta_r)^{n}$ of the external legs by $\beta'>\frac72$. This proves (i).

(ii) For a graph outside the degree-$0$ members of the order-$e^2$ $h$-linear field-free family, \Cref{rot:lem:H6census} gives strictly positive ordered degrees for every non-vanishing derived graph (the excluded vacuum case of the census is now the family itself), and \Cref{rot:cor:OF} without weakening gives the per-site bound with the factor $(1+k-r)^{p_0}$, King's ``possible addition of some power of $\ln(L^k\eps)^{-1}$'' (p.~666); the sup-norm dependence on the sources is King's (p.~664, 677), the distinguished vertex entering linearly. The two $h$-linear field-free members of order $0$ and $\lambda$ are bounded directly: the subtracted tadpole at zero coarse field is the infrared remainder $-N\sum\eta^3\delta_r^2h\,F_{r}(x,x)$ of the coincident propagator, bounded per site at unit scale, and the figure-eight minus its $\dm_1$ partner leaves a bubble of positive ordered degree times an infrared tadpole line, as in step 3 of the proof of \Cref{rot:cor:OF}. For the degree-$0$ members of the order-$e^2$ family (the exchange and seagull tadpoles with their $\dm_2$ partner) every proper ordered subgraph has $D>0$ (a proper component with the $h$-vertex is not a vacuum graph, the $h$-loop vanishes, and the source-free components are the self-energy and the photon tadpole, renormalised by the $\dm_2$ partner) and only the whole graph has $D=0$: the Taylor identity for the self-energy with its $\dm_2$ partner, Ba{\l}aban's general step at the whole graph with no final renormalisation (there is none to perform), as in King's worked case \cite[(3.81)--(3.87)]{K1}, makes every proper prefix component of degree $\ge\alpha$ and leaves the whole graph at degree $0$, so in the multiscale sum every geometric series keeps its constant except the last, $\sum_{j\le r-1}(5^j\eta)^0=r$: this is the factor $(1+r)$, the zero-degree logarithm of \Cref{rot:cor:OF}. It is necessary: the value of the family at a constant source is $u$ times the drift of \Cref{rot:rem:log}, of size $r$ at fixed $\delta_r$. The infrared counterterm lines are handled as in (i).
\end{proof}

\subsubsection*{The Cauchy estimate at the frozen wall}

\begin{proposition}[King's Cauchy estimate modulo the field-free constant]\label{rot:prop:H6I3}
Let $h=\hinf$ and $t\in\R$, and define the field-free constant $a^{(r)}_k(h)\defeq S^{(r)}_k(0,h)\big|_{A_r=\phi_r=0}$ (the whole value of the effective action at zero fields and zero curvature source) and $d_k\defeq a^{(r)}_k(h)-a^{(r+1)}_{k+1}(h)$, a number independent of $t$ and of $G$. Then on $\supp\chi_r$, with $\gamma',\beta',\sigma'$ as in \Cref{rot:prop:H6R}(i),
\[
\big|S^{(r)}_k(tG,h)-S^{(r+1)}_{k+1}(tG,h)-d_k\big|\le\eta_k|T|,\qquad\eta_k\defeq C\big[5^{-\gamma'r}\delta_r^{-\beta'}+\delta_r^{\sigma'}\big],
\]
$C$ independent of $k$ and locally uniform in $t$.
\end{proposition}
\begin{proof}
Both effective actions live on $T^{(r)}$ with the same $\chi_r$. Their Gaussian parts converge by \cite[Thm.~3.3]{K1} (the step King takes for (3.9) before the perturbative part), and their difference is a quadratic form in the fields plus a constant, the constant being part of $d_k$. In the perturbative parts, split every graph as $E(H)=E(H)|_0+E_N(H)$ when $H$ has no external scalar leg and no curvature vertex: the sum of the values $E(H)|_0$ over these graphs, at $G=0$, is by definition the field-free constant $a^{(r)}_k(h)$ minus the Gaussian constant, and the graphs with a curvature vertex and no external leg are included in $E_N$ and converge by \Cref{rot:prop:H6R}(i). So $S^{(r)}_k(tG,h)-S^{(r+1)}_{k+1}(tG,h)-d_k$ is a sum of differences $E^{(r+1)}_N(H)-E^{(r)}_N(H)$ over the finitely many graphs, each bounded by \Cref{rot:prop:H6R}(i), plus the two Taylor remainders, $\le C\delta_r^\sigma|T|$ each, plus the non-polynomial terms, included in \Cref{rot:prop:H6R}(i). This is King's \cite[(3.9)]{K1} without the field-free $h$-sector constant, which is exactly what is false in the printed form at $\int h\ne0$ (\Cref{rot:rem:log}: the order-$e^2$ $h$-linear constant drifts by $c_*\ln5\int h$ per step, so $d_k\not\to0$); $d_k$ is $t$- and $G$-independent by construction and cancels between numerator and denominator of every ratio below. Only one step, $n=1$, is used.
\end{proof}

% r03_h6w_c: proof of Theorem rot:thm:H6W, part 3: the constant-source shift on one lattice, the rate, the frozen family is Cauchy, and the assembly.

\subsubsection*{The constant-source shift on one lattice}

\begin{lemma}[the constant-source shift of the effective action]\label{rot:lem:H6SD}
Let $h,h'$ be in the family of \Cref{rot:lem:H6conv} with $h-h'=u\mathbf 1$ constant. Then on $\supp\chi_r$, for every real $t$,
\begin{equation}\label{rot:eq:H6SD}
\big|S^{(r)}_k(tG,h)-S^{(r)}_k(tG,h')\big|\le C(1+r)\,|u|\,p(\delta_r)^{n_1}\delta_r^{-3}|T|+2C\delta_r^\sigma|T|,
\end{equation}
with $C=C(t,G,h_0,|T|)$ independent of $k$ and $u$ and locally uniform in $t$, $n_1$ the exponent of \Cref{rot:prop:H6R}(ii) enlarged to absorb $(1+k-r)^{p_0}$, and $\sigma$ the remainder exponent.
\end{lemma}
\begin{proof}
\emph{Where $h$ enters.} By \cite[(3.14)--(3.17)]{K1} the Gaussian parts of $S^{(r)}_k$ are defined at the mass $m^2$ without sources, and the sources sit in $E'_r$ only, through the vertices (3.20), (3.21); the $h$-vertex is linear in $h$ with the explicit coupling $\delta_r^2\le1$. So $S^{(r)}_k(tG,h)-S^{(r)}_k(tG,h')=[P^{(r)}_k(h)-P^{(r)}_k(h')]+[R(h)-R(h')]$, $P$ the Taylor polynomial of order $\bar n=13$ and $R$ the remainder.
\emph{The remainder.} $|R(h)|,|R(h')|\le C\delta_r^\sigma|T|$ on $\supp\chi_r$, uniformly over the source family (\cite[p.~665]{K1}); this is the second term, and no derivative of $R$ in $u$ is taken.
\emph{The polynomial.} $P^{(r)}_k(h)$ is a finite sum of connected renormalised graphs with $m\ge0$ $h$-vertices, each multilinear in its source insertions; by multilinearity $P(h)-P(h')$ is the sum over the graphs with $m\ge1$ of $\sum_{i=1}^m[H$ with the $i$-th $h$-vertex carrying the constant $u$ and the others carrying $h$ or $h']$. The value of each such graph on $\supp\chi_r$ is $\sum_{\{y\},\{z\}}E^{(r)}_k(H;\{y\},\{z\})\times(\text{external fields at }\{y\})$. By \Cref{rot:prop:H6R}(ii), $|E^{(r)}_k(H)|\le C(1+r)(1+k-r)^{p_0}\delta_r^{\varrho_H}p(\delta_r)^{n_1}e^{-\delta\dist}$ with $\varrho_H\ge0$ and $C$ linear in $|u|$ for the distinguished vertex (a constant source has zero derivatives), the factor $(1+r)$ being present only for the degree-$0$ members of the order-$e^2$ $h$-linear field-free family and harmless elsewhere; the external scalar fields obey $|\phi_r(y)|\le p(\delta_r)(\lambda\delta_r)^{-1/4}$ \cite[(3.2)]{K1}, with $\delta_r^{-s/4}$ absorbed by $\varrho_H\ge s/4$ for $s$ external scalar legs, and the external photon legs are bounded likewise by (3.2), so the product of fields and kernel is $\le C|u|(1+r)(1+k-r)^{p_0}p(\delta_r)^{n_1+s}$ per position; this is the whole content of King's remark that the sources ``may be bounded when making a perturbative expansion'' and ``do not affect the leading order positivity of the action'' (p.~677). The graph is connected and its kernel decays exponentially in the tree distance of all its vertices, so the sum over all positions is $\le C|T^{(r)}|=C|T|\delta_r^{-3}$: one free position, the others tied. Collecting the finitely many graphs, with $(1+k-r)^{p_0}\le C(1+\log\delta_r^{-1})^{p_0}$ absorbed by enlarging $n_1$, gives the first term of \eqref{rot:eq:H6SD}.
\end{proof}

The bound \eqref{rot:eq:H6SD} keeps no smallness from the explicit $\delta_r^2$ of the $h$-vertex, from the running couplings or from the positive degrees: it is a per-site bound $O((1+r)p^{n_1})$ times the site count $|T|\delta_r^{-3}$. The true size of the worst term, the tree-level term $u\delta_r^2\sum\eta^3|\phi^{(r)}|^2$, is $|u|p^2\lambda^{-1/2}\delta_r^{-3/2}|T|$ on $\supp\chi_r$: a constant scalar source is a bare-mass shift, and its effect is a running-mass term of the same order as the mass term already present, which does not tend to zero. The smallness comes from the rate below, not from the effective action.

\begin{lemma}[the rate and the choice of $\vr$]\label{rot:lem:H6rate}
With $u_k=\frac12(\Delta_T(0)-\Delta_T(\eps_k))$ of \Cref{rot:lem:H6conv}, $|u_k|\le C\eps_k=C5^{-k}$ (\Cref{rot:prop:CT}(b)), and
\[
\eta'_k\defeq C(1+r)|u_k|p(\delta_r)^{n_1}\delta_r^{-3}+2C\delta_r^\sigma\le C(1+k)\,5^{-k+3\lfloor\vr k\rfloor}\big(b_0(1+k\ln5)\big)^{\mathrm pn_1}+2C\,5^{-\sigma\lfloor\vr k\rfloor}\longrightarrow0,
\]
since $\vr<\frac19<\frac13$ (\Cref{rot:lem:r}): the polynomial factor $(1+k)^{1+\mathrm pn_1}$ is beaten by the geometric $5^{-(1-3\vr)k}$. Hence, on $\supp\chi_r$, $|S^{(r)}_k(tG,h_k)-S^{(r)}_k(tG,\hinf)|\le\eta'_k|T|$, and, exponentiating under the common $\chi_r$ (same lattice, same depth, same fields),
\[
I_k(t;h_k)=e^{\theta_k(t)|T|}I_k(t;\hinf),\quad|\theta_k(t)|\le\eta'_k,\qquad\text{so}\qquad
\frac{\rho_k(t;h_k)}{\rho_k(t;\hinf)}\in[e^{-2\eta'_k|T|},e^{2\eta'_k|T|}],\qquad
|\rho_k(t;h_k)-\rho_k(t;\hinf)|\le16e^{C_Gt^2}\eta'_k|T|\to0
\]
for $k$ so large that $2\eta'_k|T|\le1$, locally uniformly in $t$.
\end{lemma}
\begin{proof}
The first display is \eqref{rot:eq:H6SD} with $u=u_k$, $\delta_r=5^{-\lfloor\vr k\rfloor}$, $1+r\le1+k$ and $p(\delta_r)\le b_0(1+k\ln5)^{\mathrm p}$; both terms tend to $0$ exactly when $0<\vr<\frac13$, and every condition of \Cref{rot:lem:r} on $\vr$ is an upper bound, so its window $\vr<\frac19$ needs no change. The second: $I_k(t;h_k)=\int\chi_re^{-S(tG,\hinf)}e^{-[S(tG,h_k)-S(tG,\hinf)]}$ lies in $[e^{-\eta'_k|T|},e^{\eta'_k|T|}]I_k(t;\hinf)$, and the same at $t=0$; divide, and use $\rho_k(t;\hinf)\le4e^{C_Gt^2}$ (\Cref{rot:lem:H6stepA}) with $e^{x}-1\le2x$ for $0\le x\le1$.
\end{proof}

\subsubsection*{The frozen family is Cauchy}

\begin{lemma}[one step at the frozen wall]\label{rot:lem:H6B}
Fix $t\in\R$ and write $r=r(k)$. For $k\ge k_0$ (so that $\delta_r<1$, $b_k,\tilde b_{k+1}\le\frac12$, $2C\delta_r^\sigma|T|\le\ln2$ and $2\eta_k|T|\le1$), with $\tilde\rho_{k+1}(t)\defeq\rho^{(r+1)}_{k+1}(t;\hinf)$ the small-field ratio of $X_{k+1}$ at depth $r+1$ (which is $r(k+1)$ or $r(k+1)+1$),
\[
\epsilon(k)\defeq|s_k(t)-s_{k+1}(t)|\le16e^{C_Gt^2}|T|\big(2C\delta_r^\sigma+\eta_k\big)+(1+4e^{C_Gt^2})\big(b_k+\tilde b_{k+1}\big),
\]
where $b_k\le2e^{(C+C_1)|T|}e^{-p(\delta_r)^2}5^{(q+q_1)k}$ and $\tilde b_{k+1}\le2e^{(C+C_1)|T|}e^{-p(\delta_r)^2}5^{(q+q_1)(k+1)}$. Consequently $\sum_k\epsilon(k)<\infty$, $(s_k(t))_k$ is Cauchy, and $W^{\rm tor}(t)\defeq\lim_ks_k(t)$ exists, locally uniformly in $t$.
\end{lemma}
\begin{proof}
\emph{(B1) The Cauchy estimate at one step.} By \Cref{rot:prop:H6I3}, on $\supp\chi_r$: $S^{(r)}_k(tG,\hinf)=S^{(r+1)}_{k+1}(tG,\hinf)+d_k+R_k(A_r,\phi_r;t)$ with $|R_k|\le\eta_k|T|$ and $d_k$ independent of the fields and of $t$. Exponentiate and integrate against the common $\chi_r$: $I^{(r)}_k(t;\hinf)=e^{-d_k}I^{(r+1)}_{k+1}(t;\hinf)e^{\theta|T|}$ with $|\theta|\le\eta_k$, and the same at $t=0$ with the same $d_k$. Dividing,
\[
\rho_k(t;\hinf)=\tilde\rho_{k+1}(t)\,e^{\theta'|T|},\quad|\theta'|\le2\eta_k,\qquad\text{so}\qquad|\rho_k(t;\hinf)-\tilde\rho_{k+1}(t)|\le\tilde\rho_{k+1}(t)\big(e^{2\eta_k|T|}-1\big).
\]
\emph{(B2) The large-field terms at the two depths.} \Cref{rot:lem:H6stepA} uses only \Cref{rot:prop:H6S1} at the given lattice and depth and \Cref{rot:prop:H6S2} for that lattice, which is depth-free; both hold for $\Xk$ at depth $r$ and for $X_{k+1}$ at depth $r+1$ ($\delta_{r+1}(X_{k+1})=\delta_r<1$), with the stated bounds on $b_k$ and $\tilde b_{k+1}$. Hence $\rho_k(t;\hinf),\tilde\rho_{k+1}(t)\le4e^{C_Gt^2}$ and, by \eqref{rot:eq:H6stepA} at both lattices,
\[
|s_k(t)-\rho_k(t;\hinf)|\le4e^{C_Gt^2}(e^{2C\delta_r^\sigma|T|}-1)+(1+4e^{C_Gt^2})b_k,\qquad
|s_{k+1}(t)-\tilde\rho_{k+1}(t)|\le4e^{C_Gt^2}(e^{2C\delta_r^\sigma|T|}-1)+(1+4e^{C_Gt^2})\tilde b_{k+1}.
\]
\emph{(B3) One step.} By (B1), (B2) and the triangle inequality,
\[
\epsilon(k)\le4e^{C_Gt^2}\big[2(e^{2C\delta_r^\sigma|T|}-1)+(e^{2\eta_k|T|}-1)\big]+(1+4e^{C_Gt^2})(b_k+\tilde b_{k+1}),
\]
and $e^x-1\le2x$ for $0\le x\le1$ gives the claim.
\emph{(B4) Summability.} With $\delta_r=5^{-\lfloor\vr k\rfloor}\le5\cdot5^{-\vr k}$ and $r=k-\lfloor\vr k\rfloor\ge(1-\vr)k$: $\delta_r^{\sigma}\le5^\sigma5^{-\sigma\vr k}$; $5^{-\gamma'r}\delta_r^{-\beta'}\le5^{\beta'}5^{-[\gamma'-(\gamma'+\beta')\vr]k}$, geometric once $\vr<\gamma'/(\gamma'+\beta')$, which holds for $\gamma'<\gamma_*$ close to $\gamma_*$ and $\beta'\in(\frac72,5)$, because the window $\vr<\gamma/(\gamma+\beta)$ of \Cref{rot:lem:r}, with $\gamma\le\gamma_*$ and $\beta\ge5$, lies strictly inside $\vr<\gamma_*/(\gamma_*+\beta')$; and $b_k+\tilde b_{k+1}\le4e^{(C+C_1)|T|}5^{q+q_1}\exp[-b_0^2(1+\lfloor\vr k\rfloor\ln5)^{2\mathrm p}+(q+q_1)k\ln5]$, a quadratic-in-$k$ negative exponent against a linear positive one ($\mathrm p\ge1$, \Cref{rot:lem:r}), hence $\le Ce^{-ck^2+C'k}$. So $\sum_{k\ge k_0}\epsilon(k)<\infty$, and for $k_0\le k<k'$, $|s_k(t)-s_{k'}(t)|\le\sum_{j=k}^{k'-1}\epsilon(j)\le\sum_{j\ge k}\epsilon(j)\to0$. All constants are uniform for $t$ in compacts. No uniformity in the number of steps is used or claimed: a sequence with summable consecutive differences is Cauchy.
\end{proof}

This lemma alone is King's torus theorem \cite[Thm.~2.1]{K1} with the frozen wall $\hinf$ as his scalar source, read on the normalised ratio: the differences from King's argument \cite[(3.10)--(3.13)]{K1} are that the drift-free stability bound $|\ln Z|\le C|T|$ \cite[(2.23)]{K1}, false at a wall, is replaced by the polynomial envelope of \Cref{rot:prop:H6S1,rot:prop:H6S2}, that the field-free $h$-sector constant, which does not converge, is removed by the ratio, and that the uniformity in the number of steps is replaced by summation over single steps.

\subsubsection*{Assembly}

\begin{proof}[Proof of \Cref{rot:thm:H6W}]
Fix $t\in\R$. By \Cref{rot:lem:H6stepA} at the two walls and \Cref{rot:lem:H6rate},
\[
|r_k(t)-s_k(t)|\le|r_k(t)-\rho_k(t;h_k)|+|\rho_k(t;h_k)-\rho_k(t;\hinf)|+|\rho_k(t;\hinf)-s_k(t)|\longrightarrow0 ,
\]
and by \Cref{rot:lem:H6B} $s_k(t)\to W^{\rm tor}(t)$. Hence $\lim_kr_k(t)=\lim_k\langle e^{tF(G)}\rangle_{\fc,\Xk}$ exists and equals $W^{\rm tor}(t)$, for every real $t$; by \Cref{rot:lem:H6conv} this is the limit of the torus-convention ratios at the moving wall $h_k=h_0-\Delta_T(\eps_k)/2$, which coincide with $r_k(t)$ term by term, and of those at the frozen wall $\hinf=h_0-\Delta_T(0)/2$. By \Cref{rot:lem:H6A} the convergence is locally uniform on $\C$, the limit $\Wv_\fc(zG)$ is entire, and all moments of $F(G)$ converge. The limits at $h_0$ itself in the two conventions are not claimed equal, and nothing below uses that.

\emph{What the proof does not use.} No value of the constants $d_k$; no Wick scheme; no convergence of $\Ztor_k(t;h)$ or $\Ztor_k(0;h)$ separately (both drift like $\eps_k^{\pm q}$); no comparison of two lattices, no rotated lattice, no skew torus, no hybrid flow and no centred blocks; no change of mass; the bridge \Cref{rot:prop:CT} only through the exact identity (c), the boundedness of $\Delta_T$ and the rate $|\Delta_T(s)-\Delta_T(0)|\le Cs$ of (b); and the sign of the wall nowhere, when the lower bound is taken from \Cref{rot:thm:FS}. The statement is at fixed $\fc$, $G$ and $t$, with $\eps_k\to0$ along $\eps_k=5^{-k}$; nothing depends on the later infimum over configurations or on the wall-height index of the shell construction.
\end{proof}


\subsubsection*{The forward bound and the assembly}

\begin{proposition}[the forward bound]\label{rot:prop:Kforward}
For $T\in\{T_1,T_2\}$ and every real source $s=(zG,h)$ with $h$ a wall $-n\psi$ or its transport,
\[
\Zun_{\Yk,T}(s)\le(1+\delta_k)\,e^{-c_k\sS_{\Xk}(h)}\,\Zun_{\Xk,T}(s),\qquad\delta_k=\delta_k(T,s)\to0,
\]
with the constant $c_k$ of \Cref{rot:lem:C}, independent of $T$ and of the source, and $|c_k|\le C\eps_k^{-q}$. No bound on $c_k$ is used.
\end{proposition}
\begin{proof}
Write $I_Y\defeq\int\chi_re^{-S^{(r)}}$ and $I_X\defeq\int\chi_re^{-\hs^{(r)}}$. By \Cref{rot:prop:H3}, $\Zun_Y\le e^{C\delta_r^\sigma|T|}I_Y+B$ with $B\defeq e^{-p(\delta_r)^2+C|T|}\eps^{-q}$. By \Cref{rot:prop:H5} (at the wall, \Cref{rot:thm:FS}), $B\le b\,\Zun_Y$ with $b\defeq e^{-p(\delta_r)^2+(C+C_1)|T|}\eps^{-q-q'}$; since $p(\delta_r)^2$ grows at least quadratically in $k$ (\Cref{rot:lem:r}) while $\ln\eps^{-q-q'}=(q+q')k\ln5$ grows linearly, $b\to0$, so $\Zun_Y\le(1-b)^{-1}e^{C\delta_r^\sigma|T|}I_Y$: the large-field term is absorbed on the left, which is why no bound on $c_k$ and no lower bound for the axis lattice is needed. On $\supp\chi_r$, \Cref{rot:thm:Asrc} gives $e^{-S}\le e^{-\hs-c_k\sS_{\Xk}(h)+\vartheta_k|T|}$, so $I_Y\le e^{-c_k\sS_{\Xk}(h)+\vartheta_k|T|}I_X$. By \Cref{rot:prop:H4}, $I_X\le e^{C(1+k)^p\delta_r^\sigma|T|}\Zun_X$. So $\delta_k=(1-b)^{-1}e^{(2C(1+k)^p\delta_r^\sigma+\vartheta_k)|T|}-1$, which tends to $0$ by \Cref{rot:lem:r}. Multiplying $\Zun$ and the small-field integrals of either lattice by a positive constant changes nothing above, which is the $c$-number blindness used in \Cref{rot:prop:CT}(b).
\end{proof}

\begin{proposition}[mass-increasing walls only]\label{rot:prop:Konly}
Let $\cS^-\defeq\{(tG,h_0):t\in\R\}$ with $h_0=-n\psi$, $n>0$, a family not closed under $h\mapsto-h$, and assume the forward bounds of \Cref{rot:prop:Kforward} on $T_1$ for $s\in\cS^-$ and on $T_2$ for their $S$-images, with the common constant $c_k$. Then the conclusions \textup{(a)} and \textup{(c)} of \Cref{rot:thm:sandwich} hold for the sources of $\cS^-$, with $v=v'=0$, $c=c'=c_k$ and $c_k\to0$ (\Cref{rot:prop:CK}); the sources $(tG,+n\psi)$ do not occur.
\end{proposition}
\begin{proof}
Part (a) of \Cref{rot:thm:sandwich} uses only the two forward bounds, one source at a time: the upper bound is the forward bound on $T_1$, the lower bound the forward bound on $T_2$ for $S\cdot s$ pulled back by \Cref{rot:lem:inv}, with $\sS_{\Lambda}(h)=\int h+O(\eps^N)$ (\Cref{rot:lem:lsum}) and $N>q$. Part (b) is the only place where the pair $(G,\pm h)$ enters, and its conclusion $c+c'\to0$ is here supplied directly by $c=c'=c_k\to0$; the interval of (a) being non-empty gives $2c_k\int h_0\le o(1)$ with $\int h_0<0$, that is $\liminf c_k\ge0$, consistently. Part (c) then follows as in that proof: both ends of the interval in (a) tend to $1$.
\end{proof}

\begin{proof}[Proof of \Cref{rot:thm:CT}]
Fix $\fc=(J,\psi,n)$ in the class and $G$, and let $\cS^-=\{(tG,h_0):t\in\R\}$, $h_0=-n\psi$.
\emph{Step 1 (comparison of the effective actions).} \Cref{rot:thm:Asrc} on $T_1$ for $s\in\cS^-$ and on $T_2$ for $S\cdot s$ (\Cref{rot:lem:walledalg}(b) transports the own-frame sources; the $S$-image of a wall is a wall on $T_2$), with the common constant $c_k$ of \Cref{rot:lem:C}, which does not depend on the torus. Inputs: \Cref{rot:thm:H1S}, \Cref{rot:prop:BD}, \cite[Thm.~3.3]{K1}, \Cref{rot:conv:C}(ii), \Cref{rot:rem:CMrg}.
\emph{Step 2 (forward bounds).} \Cref{rot:prop:Kforward} on $T_1$ and on $T_2$ for these sources. Inputs: \Cref{rot:prop:H3,rot:prop:H4,rot:prop:H5} on both tori, hence \Cref{thm:a06:BALA}, \Cref{rot:prop:CT} and \Cref{rot:conv:C} throughout; on $T_1$ every $Y$-frame step is a skew step (\Cref{rot:prop:SK}); on $T_2$ only the first hybrid step is skew, with image terms $O(e^{-c\,50MJ/\eps})$. The instance of \Cref{rot:prop:H5} at the wall is \Cref{rot:thm:FS}.
\emph{Step 3 (sandwich).} \Cref{rot:prop:Konly} with \Cref{rot:prop:CK}: conclusion (c) of \Cref{rot:thm:sandwich} for the balanced pair $s_1=(tG,h_0)$, $s_0=(0,h_0)$, whose scalar sources have the same integral, gives
\[
u_k(t)\defeq\langle e^{tF(G)}\rangle_{\fc,\Yk}-\langle e^{tF(G)}\rangle_{\fc,\Xk}\to0\qquad\text{for every real }t .
\]
\emph{Step 4 (complex coefficient, moments).} By \Cref{lem:s02:GD} and \Cref{rot:lem:CM}, $1\le\langle e^{tF}\rangle\le e^{t^2q_\Lambda/2}$ with $q_\Lambda\to q$ on both lattices (\Cref{rot:lem:q}), and $|\langle e^{zF}\rangle|\le\langle e^{(\Re z)F}\rangle$; so $u_k$ is bounded on strips uniformly in $k$ and tends to $0$ on $\R$; Vitali's theorem \cite{Hor90} gives (i), and Cauchy's estimates the moments.
\emph{Step 5 (existence).} By \Cref{rot:thm:H6W} the axis family converges for real $t$; by the same argument it converges locally uniformly on $\C$ to an entire function $\Wv_\fc(zG)$, and by (i) so does the rotated family, to the same limit. This is (ii).
\emph{Step 6 (other generators).} \Cref{rot:lem:K}(c): the statement for $(R^{-1},T_1)$ and for the other planes follows from the one for $(R,T_1)$ by the passive reflections and permutations, which preserve $\Xk$ and $T_1$ and transport own-frame sources to own-frame sources, and the class of configurations is invariant under them. This is (iii).
No reverse step (``switching around'', \cite[p.~337]{K2}), no comparison of two spacings, no Wick scheme and no value of any $h$-linear constant, no convergence of the rotated family on its own and no infinite-volume limit is used: the two-lattice content is Steps~1--4, and existence is the one-lattice question of Step~5.
\end{proof}

\begin{proof}[Completion of the proof of \Cref{rot:hyp:K}]
By \Cref{rot:lem:K}(c) it suffices to treat $R=R_{12}(\thz)$, and by \Cref{rot:rem:class} the configurations to be treated are those of the class of \Cref{rot:thm:CT}. For such $\fc$ and $G$, Steps~1--2 of the proof of \Cref{rot:thm:CT} give the bound \hypF\ of \Cref{rot:hyp:F} on both tori for the sources $(tG,h_0)$, in total form: $v_k(T_1)=v_k(T_2)=0$ and $c_k(T_1)=c_k(T_2)=c_k$, with $|c_k|\le C\eps_k^{-q}$ (\Cref{rot:lem:C}), so that \hypFz\ holds for these sources, \hypV\ and \hypCc\ hold trivially, and moreover $c_k\to0$ (\Cref{rot:prop:CK}). \Cref{rot:prop:Konly} supplies conclusion (c) of \Cref{rot:thm:sandwich} for the family $\cS^-$ without the sources $(tG,+n\psi)$, and \Cref{rot:cor:walled} then gives $\langle e^{zF(G)}\rangle_{\fc,R^{-1}\Lk}-\langle e^{zF(G)}\rangle_{\fc,\Lk}\to0$ locally uniformly in $z$; with the existence of $\lim_k\langle e^{tF(G)}\rangle_{\fc,\Lk}$ for real $t$ (\Cref{rot:thm:H6W}) its last clause gives \Kz\ for $\fc$ and $G$.
\end{proof}
